% rrxiv-template.tex — skeleton paper demonstrating every rrxiv.cls feature.
%
% Compile with:
%   tectonic --keep-intermediates rrxiv-template.tex
%   # or: pdflatex rrxiv-template.tex && pdflatex rrxiv-template.tex
%
% Replace the metadata, authors, abstract, and content with your own.
% Keep at least one claim, one evidence block, and one citation so the
% template paper itself parses to a non-trivial CIR.

\documentclass{rrxiv}

% ---- Required rrxiv metadata ----
\rrxivid{my-paper-0001}
\rrxivversion{v1}
\rrxivprotocolversion{0.1.0}
\rrxivlicense{CC-BY-4.0}
\rrxivtopics{example-topic,scientific-publishing}

% ---- Standard LaTeX metadata ----
\title{Title of the paper, as a single falsifiable claim if possible}

\author[1]{First Author}
\author[2]{Second Author}

\affil[1]{Department of X, University Y}
\affil[2]{Independent}

\date{2026-05-04}

\begin{document}
\maketitle

\begin{abstract}
A short, plain-text abstract. Math can be inline like $E = mc^2$, but
pretend the parser ignores math when extracting the abstract for the CIR.
The abstract is one of the few non-environment fields the parser cares
about, so keep it standalone-readable.
\end{abstract}

% ============================================================
\section{Introduction}
\label{sec:intro}
% ============================================================

Open with the motivation in a few paragraphs. State what's new and why a
reader should care. Cite prior work: \cite{example-prior-work}.

\begin{rrxivremark}[Optional remarks]
\label{rem:context}
Use \texttt{rrxivremark} for asides, methodological notes, or design
context that don't constitute claims. The parser captures them as
\texttt{remark} nodes, separate from the claim graph.
\end{rrxivremark}

% ============================================================
\section{Setup and definitions}
\label{sec:setup}
% ============================================================

\begin{scope}[Conditions under which the claims below hold]
\label{scope:main}
Use \texttt{scope} environments to declare the regime, models, datasets,
or assumptions your claims depend on. The parser maps the body to the
\texttt{scope} field of every claim emitted between this scope and the
next one.
\end{scope}

\begin{observation}[A factual observation, no claim attached]
\label{obs:setup}
Observations are factual or empirical statements you make in passing
without asserting them as the paper's main contribution. They get a
node in the CIR but are typed differently from claims.
\end{observation}

% ============================================================
\section{Main results}
\label{sec:results}
% ============================================================

\begin{claim}[A single load-bearing assertion]
\label{claim:main}
State the claim as a single falsifiable assertion in plain language.
A reader should be able to read this paragraph in isolation and know
what is being asserted, with no reference to surrounding context.
\end{claim}

\begin{evidence}[Evidence for claim:main]
\label{ev:main}
Sketch or describe the evidence that supports the claim above. The
parser pairs the most recent \texttt{claim} with the next
\texttt{evidence} block by default; cross-paper supports relations
should use the \texttt{\textbackslash supports} edge below.
\end{evidence}

% Inline edge: this claim depends on a prior claim in another paper.
\dependson{my-paper-0001:claim:main}{rrxiv-0001:claim:queryability}

% Inline edge: this claim supports a follow-on claim.
\supports{my-paper-0001:claim:main}{my-paper-0001:claim:secondary}

\begin{claim}[A secondary claim that builds on the main one]
\label{claim:secondary}
A second claim. Note the inline \texttt{\textbackslash supports} edge
above declares this claim is supported by \texttt{claim:main}.
\end{claim}

% ============================================================
\section{Discussion}
\label{sec:discussion}
% ============================================================

\begin{openquestion}[Something this paper does not resolve]
\label{oq:scaling}
Open questions are first-class in the CIR. Use them when you want to
flag what your paper does not address but a future paper or annotation
might.
\end{openquestion}

% Inline edge: this paper extends a claim from another paper.
\extendsclaim{my-paper-0001:claim:main}{some-foundational-paper:claim:base-result}

% Inline edge: this paper contradicts a claim from another paper.
% Use sparingly; the parser flags contradictions for community review.
% \contradicts{my-paper-0001:claim:main}{some-paper:claim:opposite-result}

% ============================================================
\section{Conclusion}
\label{sec:conclusion}
% ============================================================

A paragraph or two summarising the contribution. Restate the main claim
and where it fits in the literature.

% ============================================================
% Bibliography
% ============================================================

\bibliographystyle{plainnat}
\bibliography{rrxiv-template}

\end{document}
